% Suggested methods text, not experimental results. Adapt axis/frame labels to
% the paper and report actual acquisition/filter settings after validation.
\paragraph{Breakaway assistance evaluation.}
We evaluate breakaway assistance using repeated, directed rest-to-motion
trials with matched initial configuration, grip point, and distal payload.
Let $u_i(t)$ denote the measured interaction force (translation) or moment
(rotation) along the instructed axis, expressed at the common tool reference
point. The signal is corrected for sensor bias and distal-body gravity;
dynamic correction is applied only when the required payload parameters are
available. With $t_p$ the detected effort onset and $t_m$ the first sustained
directed motion, we report
\begin{equation}
 U_{\mathrm{break},i}=\max_{t\in[t_p,t_m]}|u_i(t)-\bar u_{i,\mathrm{rest}}|,
 \qquad T_{\mathrm{onset},i}=t_m-t_p.
\end{equation}
Trials without a valid onset are reported as failures, not assigned zero effort.
We also report effort RMS, unintended rotation during translation, and command
saturation frequency. The shared friction scale multiplies the breakaway term;
therefore its zero setting disables assistance irrespective of the nominal
breakaway fraction. Decay-speed settings at zero fraction are physically
equivalent and are not interpreted as separate treatment effects.

\paragraph{Damping evaluation.}
For a prescribed release time $t_r$, we evaluate post-release translational
and angular path lengths,
\begin{equation}
 D_{\mathrm{release}}=\int_{t_r}^{T}\|v(t)\|\,dt,\qquad
 A_{\mathrm{release}}=\int_{t_r}^{T}\|\omega(t)\|\,dt,
\end{equation}
and the time required to remain below predefined linear and angular speed
thresholds for a sustained interval. Release validity is checked using measured
interaction effort and trace inspection. Trials with continued hand contact
are not used to estimate free-decay behavior. Interaction work is
\begin{equation}
 E_h=\int_0^T W_h(t)^T V(t)\,dt,
\end{equation}
with power-conjugate wrench and twist expressed in the same frame and at the
same reference point. These are empirical damping and interaction-energy
measurements, not a proof of passivity; negative port work may include return
of stored energy. The damping sweep independently switches translational and
rotational damping off/on in a $2\times2$ design. Conditions that cannot be
completed safely are reported explicitly, not silently omitted from the comparison.

\paragraph{Local apparent inertia.}
Over short windows with approximately constant local dynamics and limited
off-axis motion, we fit the integral model
\begin{equation}
 \int_{t_a}^{t_b}u_i(t)\,dt =
 M_{\mathrm{app},i}\Delta v_i+B_i\Delta x_i+
 C_i\int_{t_a}^{t_b}\tanh(v_i/v_0)\,dt+b_i(t_b-t_a).
\end{equation}
For rotational axes, the corresponding angular quantities are used and the
inertial coefficient has units of $\mathrm{kg\,m^2}$. We report the fit quality,
excitation/conditioning checks, timing assumptions, and saturation activity.
The local scalar fit is not an identification of the complete coupled
operational-space inertia matrix. Repeated trials, rather than individual
sensor samples, form the statistical replication unit.

\paragraph{Rotation leakage within the ablations.}
For directed translation along $i\in\{x,y,z\}$ we compute
\begin{equation}
 L_i=\frac{180}{\pi}
 \frac{\int\|\omega(t)\|\,dt}{\int |v_i(t)|\,dt}
 \quad[\mathrm{deg/m}],
\end{equation}
over the instructed motion interval. Trials below a predefined minimum
translation path are excluded from this ratio and reported separately.
This metric is collected within each parameter sweep, rather than in a
separate rotation-leakage experiment.
